%%%%%%%%%%%%%%%%%%%%%%%%%%%%%%%%%%%%%%%%%%%%%%%%%%
%%                                              %%
%%        Comandos usados no documento          %%
%%                                              %%
%%%%%%%%%%%%%%%%%%%%%%%%%%%%%%%%%%%%%%%%%%%%%%%%%%
% Capa e Folha de Rosto
\renewcommand{\imprimircapa}{
%%%%%%%%%%%%%%%%%%%%%%%%%%%%%%%%%%%%%%%%%%%%%%%%%%
%%                                              %%
%%              Início da Capa                  %%
%%                                              %%
%%%%%%%%%%%%%%%%%%%%%%%%%%%%%%%%%%%%%%%%%%%%%%%%%%

\thispagestyle{empty}

\begin{figure}
    \centering
    \includegraphics[]{lib/Logoufrn.jpg}
\end{figure}

\begin{center}
    \large{\textbf{\imprimirinstituicao}}

    \vspace{\stretch{2}}
    
    \large{\MakeUppercase{\textbf{\imprimirautor}}}
    
    \vspace{\stretch{2}}
    
    \large{\MakeUppercase{\textbf{\imprimirtitulo \imprimirsubtitulo}}}
    
    \vspace{\stretch{3}}
    
    \large{\textbf{\imprimirlocal}} \\
    \large{\textbf{\imprimirdata}}
\end{center}

\newpage
%%%%%%%%%%%%%%%%%%%%%%%%%%%%%%%%%%%%%%%%%%%%%%%%%%
%%                                              %%
%%                  Fim da Capa                 %%
%%                                              %%
%%%%%%%%%%%%%%%%%%%%%%%%%%%%%%%%%%%%%%%%%%%%%%%%%%
}
\renewcommand{\imprimirfolhaderosto}{
%%%%%%%%%%%%%%%%%%%%%%%%%%%%%%%%%%%%%%%%%%%%%%%%%%
%%                                              %%
%%          Início da Folha de Rosto            %%
%%                                              %%
%%%%%%%%%%%%%%%%%%%%%%%%%%%%%%%%%%%%%%%%%%%%%%%%%%

\thispagestyle{empty}

\begin{center}
    % Imprime o nome do autor depois coloca um espaçamento vertical de 3 vezes o espaçamento existente entre os textos
    \large{\MakeUppercase{\imprimirautor}}
    \vspace*{\stretch{3}}
    
    % Imprime o título do trabalho depois coloca um espaçamento vertical padrão dentre os textos da página
    \large{\MakeUppercase{\imprimirtitulo: \imprimirsubtitulo}}
    \vspace{\stretch{1}}
    
    % Inicia a mini página alinhada horizontalmente com o comprimento da linha sendo a metade da mesma imprimindo o preâmbulo, o orientador do trabalho e o co-orientador (se houver). Após a mini página o espaçamento vertical
    \hfill
    \begin{minipage}{.5\linewidth}
        \small\imprimirpreambulo
        \\ [0.6cm]
        \small \imprimirorientadorRotulo \imprimirorientador.\\
        \small \imprimircoorientadorRotulo
        \small \imprimircoorientador.
    \end{minipage}
    \vspace{\stretch{1}}
    
    % Imprime local e data do trabalho
    \large{\imprimirlocal}\\
    \large{\imprimirdata}
\end{center}

\newpage
%%%%%%%%%%%%%%%%%%%%%%%%%%%%%%%%%%%%%%%%%%%%%%%%%%
%%                                              %%
%%            Fim da Folha de Rosto             %%
%%                                              %%
%%%%%%%%%%%%%%%%%%%%%%%%%%%%%%%%%%%%%%%%%%%%%%%%%%
}

% subtítulo
\providecommand{\imprimirsubtitulo}{}
\newcommand{\subtitulo}[1]{\renewcommand{\imprimirsubtitulo}{#1}}

% criar fonte em imagens, tabelas e quadros
\newcommand{\source}[1]{\legend{\textbf{Fonte:} {#1}}}    

% texto em arial
%\renewcommand{\sfdefault}{phv}
%\renewcommand{\rmdefault}{phv}

% texto em Times
%\renewcommand{\sfdefault}{ptm}
%\renewcommand{\rmdefault}{ptm}

%%%%%%%%%%%%%%%%%%%%%%%%%%%%%%%%%%%%%%%%%%%%%%%%
%%          Informações do documento          %%
%%%%%%%%%%%%%%%%%%%%%%%%%%%%%%%%%%%%%%%%%%%%%%%%

\titulo{TITULO DO TRABALHO}
\tipotrabalho{Monografia (bacharel em Sistema de Informação) }
\autor{AUTOR DO TRABALHO}
\orientador{ORIENTADOR DO TRABALHO}     % nível acadêmico + nome
\local{Caicó- RN}
\data{\the\year}
\instituicao{
    UNIVERSIDADE FEDERAL DO RIO GRANDE DO NORTE\par
    CENTRO DE ENSINO SUPERIOR DO SERIDÓ \par
    DEPARTAMENTO DE COMPUTAÇÃO E TECNOLOGIA \par
    BACHARELADO EM SISTEMAS DE INFORMAÇÃO \par
}


% Rótulo do orientador
\renewcommand{\imprimirorientadorRotulo}{Orientador(a): }
\renewcommand{\imprimircoorientadorRotulo}{Co-orientador(a): }

\preambulo{Trabalho de conclusão de curso apresentado ao curso de graduação em Sistemas de Informação, como parte dos requisitos para obtenção do título de Bacharel em Sistemas de Informação da Universidade Federal do Rio Grande do Norte.}

%%%%%%%%%%%%%%%%%%%%%%%%%%%%%%%%%%%%%%%%%%%%%%%%%%
%%                                              %%
%%        Comandos usados no documento          %%
%%                                              %%
%%%%%%%%%%%%%%%%%%%%%%%%%%%%%%%%%%%%%%%%%%%%%%%%%%

%%%%%%%%%%%%%%%%%%%%%%%%%%%%%%%%%%
%%                              %%
%%            Quadros           %%
%%                              %%
%%%%%%%%%%%%%%%%%%%%%%%%%%%%%%%%%%
\newcommand{\quadroname}{Quadro}

\newfloat[chapter]{quadro}{loq}{\quadroname}
\newlistof{listofquadros}{loq}{\listofquadrosname}
\newlistentry{quadro}{loq}{0}

% configurações para atender às regras da ABNT
\setfloatadjustment{quadro}{\centering}
\counterwithout{quadro}{chapter}
\renewcommand{\cftquadroname}{\quadroname\space} 
\renewcommand*{\cftquadroaftersnum}{\hfill--\hfill}

\setfloatlocations{quadro}{hbtp}


%%%%%%%%%%%%%%%%%%%%%%%%%%%%%%%%%%%%%%%%%%%%%%%%%%
%%                                              %%
%%   Configurações de aparência do PDF final    %%
%%                                              %%
%%%%%%%%%%%%%%%%%%%%%%%%%%%%%%%%%%%%%%%%%%%%%%%%%%

%%%%%%%%%%%%%%%%%%%%%%%%%%%%%%%%%%
%%                              %%
%%      Informações do PDF      %%
%%                              %%
%%%%%%%%%%%%%%%%%%%%%%%%%%%%%%%%%%
\makeatletter
\hypersetup{
    %  	pagebackref=true,
		pdftitle={\@title}, 
		pdfauthor={\@author},
        pdfsubject={\imprimirpreambulo},
        pdfcreator={LaTeX with abnTeX2},
		pdfkeywords={abnt}{latex}{abntex}{abntex2}{trabalho acadêmico}, 
    %   false: boxed links; true: colored links
		colorlinks=true,
    % 	color of internal links
        linkcolor=blue,
    %   color of links to bibliography
        citecolor=blue,
    %   color of file links
        filecolor=magenta,
		urlcolor=blue,
		bookmarksdepth=4
}
\makeatother

%%%%%%%%%%%%%%%%%%%%%%%%%%%%%%%%%%
%%                              %%
%%         Indentação           %%
%%                              %%
%%%%%%%%%%%%%%%%%%%%%%%%%%%%%%%%%%

% O tamanho do parágrafo é dado por:
\setlength{\parindent}{1.5cm}

% Controle do espaçamento entre um parágrafo e outro:
\setlength{\parskip}{0.2cm}  % tente também \onelineskip


%%%%%%%%%%%%%%%%%%%%%%%%%%%%%%%%%%
%%                              %%
%%         Algoritmos           %%
%%                              %%
%%%%%%%%%%%%%%%%%%%%%%%%%%%%%%%%%%

% Algoritmos
% \floatname{algorithm}{Algoritmo}%Algoritmo
% \renewcommand{\listalgorithmname}{LISTA DE ALGORITMOS}

\renewcommand{\algorithmicindent}{3.0em}
\renewcommand{\algorithmicrequire}{\textbf{entrada}}
\renewcommand{\algorithmicensure}{\textbf{garanta}}
\renewcommand{\algorithmicreturn}{\textbf{retorne}}
\renewcommand{\algorithmicor}{\textbf{ou}}
\renewcommand{\algorithmicand}{\textbf{e}}
\renewcommand{\algorithmicend}{\textbf{fim}}
\renewcommand{\algorithmicif}{\textbf{se}}
\renewcommand{\algorithmicelse}{\textbf{sen\~ao}}
\renewcommand{\algorithmicthen}{\textbf{ent\~ao}}
\renewcommand{\algorithmicfor}{\textbf{para}}
\renewcommand{\algorithmicforall}{\textbf{para cada}}
\renewcommand{\algorithmicrepeat}{\textbf{repita}}
\renewcommand{\algorithmicuntil}{\textbf{até que}}
\renewcommand{\algorithmicwhile}{\textbf{enquanto}}
\renewcommand{\algorithmicdo}{\textbf{fa\c{c}a}}

